% !PRINTED 683
\weberpage{683}
%<<CONT
and if we therefore assume, which is plainly no restriction of
generality, that $\xi_1$ is a non-vanishing proper fraction, then this
formula contradicts the assumption that

\[
\varepsilon_1, \varepsilon_2, \ldots, \varepsilon_{\nu-1}
\]

is a fundamental system, because the determinant $L_{\nu-1}$ would be
 diminished if $\varepsilon_1$ were replaced by $\varepsilon_0$.

If we have a fundamental system, we may take an arbitrary
 system of integral rational numbers $\xi_1, \xi_2, \ldots, \xi_{\nu-1}$ as
system of exponents and form a unit $\varepsilon$ with these exponents:

\[
\varepsilon = \varepsilon_1^{\xi_1} \varepsilon_2^{\xi_2} \ldots \varepsilon_{\nu-1}^{\xi_{\nu-1}}.
\]

There thus remains the question to be answered, how far
a unit is determined by its system of exponents.

Let us assume that $\varepsilon', \varepsilon''$ are two units with the same
system of exponents; then the system of exponents of the
unit $\varepsilon' : \varepsilon'' = \varrho$ consists of nothing but zeros, and consequently the conjugate
 logarithms of $\varrho$ are all equal to zero. Hence $\varrho$ is an
integer with the property that the absolute values of all
numbers conjugate to $\varrho$ are equal to 1, and therefore, as was proved in
\S.\,175, 2., a \gesperrt{root of unity}. With this the
following theorem is proved:

I. \gesperrt{In the field $\Omega$ there is a system of $\nu - 1$ fundamental
 units}

\weq{5}{\varepsilon_1, \varepsilon_2, \ldots, \varepsilon_{\nu-1},}

\gesperrt{which has the property that in the form}

\weq{6}{\varepsilon = \varrho\, \varepsilon_1^{\xi_1} \varepsilon_2^{\xi_2} \ldots \varepsilon_{\nu-1}^{\xi_{\nu-1}}}

\gesperrt{all units of the field are contained, each only once,
when $\xi_1, \xi_2, \ldots, \xi_{\nu-1}$ run through all integral rational numbers and $\varrho$ through all
roots of unity present in $\Omega$.}

It is now easy to derive all the other fundamental systems
from one of them, by putting in (6) for the exponents $\nu - 1$ different
 systems of integers $\xi_{i,k}$ whose determinant
$= \pm 1$. For if one sets

\[
\varepsilon_i' = \varrho_i\, \varepsilon_1^{\xi_{1,i}} \varepsilon_2^{\xi_{2,i}} \ldots \varepsilon_{\nu-1}^{\xi_{\nu-1,i}},
\]

then it follows from (3):

\[
L(\varepsilon_1', \varepsilon_2', \ldots, \varepsilon_{\nu-1}') = \pm L(\varepsilon_1, \varepsilon_2, \ldots, \varepsilon_{\nu-1}).
\]
